% sections/related.tex -- Section 2 (paper_v2, cut in C3 from the full version's related work).
% Every work cited was read for this project; each claim was checked against the paper's text.
% B1 (25 November) brought in the author's review of the Vietnamese related work (paper_v2_vn,
% L10): the paragraph on evaluating search methods, which frames this study as search against
% manual design with random draws as the control, and the sentence on time-ignoring splits moved to
% the conclusion's limitations.
\section{Related work}
\label{sec:related}

\paragraph{Meta-path-based recommendation.} MCRec weighs a meta-path-based context of each
user--item pair with co-attention \citep{hu2018leveraging}. HGRec \citep{shi2020hgrec},
FHRec \citep{mu2021fhrec} and AMERec \citep{yan2021attention} learn one embedding per meta-path and
fuse the embeddings with attention as HAN \citep{wang2019heterogeneous} does, and metapath2vec learns node
embeddings from random walks that follow a meta-path \citep{dong2017metapath2vec}. All of these
take their meta-paths as given.

\paragraph{Automatic meta-path discovery.} Query-based methods find the $k$ best meta-paths
connecting two given nodes \citep{zhu2020topk}. Search-based methods learn which paths (or richer
structures) a model should use. GEMS evolves sets of meta-structures with a genetic algorithm
\citep{han2020genetic} and RMS extends meta-path sets with reinforcement learning
\citep{ning2022automatic}. Both judge candidates by the recommender itself. Training a recommender
per candidate is slow: in DiffMG's comparison, GEMS's search took \val{diffmg-gems-min} to
\val{diffmg-gems-max} GPU minutes per dataset, against at most \val{diffmg-own-max} for three runs of
DiffMG's search \citep[Table~5]{ding2021diffmg}. One-shot searches in the style of DARTS
\citep{liu2019darts} cut this cost by training only once: all candidates are trained together in
one supernet. DiffMG searches for a meta-graph that guides a graph neural network for node
classification and user--item link prediction, then keeps the best of three search runs by
validation score \citep{ding2021diffmg}. LMSPS samples paths progressively from a large candidate
set \citep{li2024longrange}. We adapt LMSPS to link prediction. Unlike GEMS and RMS, it judges its
sets with its own network, a stand-in for HAN (Sections~\ref{sec:res-weights}
and~\ref{sec:conclusion}).

\paragraph{Evaluating search methods.} In HAN, LMSPS compares its paths with HAN's original paths,
which were chosen by hand \citep[Table~4]{li2024longrange}. It compares them with random paths only
in its own network, where the random ones score lower \citep[Table~1]{li2024longrange}. A random
control is needed: in neural architecture search, random search with early stopping did at least as well as ENAS (a
leading search method) on two standard benchmarks \citep{li2019random}. On average,
state-of-the-art search policies ``perform similarly to the random policy''
\citep{yu2020evaluating}. Among meta-path searches, RMS beats a random control: in the same time,
the control trains the recommender on random sets and keeps the one with the best validation score
\citep{ning2022automatic}. GEMS and DiffMG do not compare with random structures
\citep{han2020genetic,ding2021diffmg}. We compare the searched sets with the manual ones in HAN and
metapath2vec (where the paths are used), with random draws from the search's own candidates as the
control.

\paragraph{Evaluating recommenders.} Simple baselines often hold their own: of
\val{dacrema-algos} recent neural recommenders, only \val{dacrema-repro} could be reproduced with
reasonable effort, and \val{dacrema-beaten} of those could often be beaten by simple heuristics
\citep{dacrema2019progress}. We therefore report popularity and a type-blind GAT, a strong baseline
on heterogeneous graphs \citep{lv2021hgb}. Sampled metrics can order models differently from a
full ranking \citep{krichene2022sampled}, so we report every path set on full ranking as well.
